\chapter{Milestone 2: Classifiers Accuracy}

\section{Obiettivo della Milestone 2}

La \textbf{Milestone 2} ha come obiettivo principale il confronto dell'accuratezza di tre classificatori di Machine Learning applicati al dataset costruito nella Milestone 1.

L'input della Milestone 2 è quindi il file CSV prodotto nella Milestone 1, contenente le metriche delle classi e la relativa etichetta di \textbf{bugginess}.

Le metriche da confrontare non si limitano alla semplice accuracy, ma includono:

\begin{itemize}
    \item \textbf{Precision};
    \item \textbf{Recall};
    \item \textbf{AUC};
    \item \textbf{Kappa};
    \item \textbf{NPofB20}.
\end{itemize}

\begin{notaesame}
La Milestone 2 non richiede solo di eseguire tre classificatori, ma di confrontarli criticamente usando più metriche. Il miglior classificatore può cambiare a seconda della metrica considerata.
\end{notaesame}

\section{Classificatori da confrontare}

\begin{notaslide}
I tre classificatori richiesti per l'esperimento sono \textbf{RandomForest}, \textbf{NaiveBayes} e \textbf{IBk}.
\end{notaslide}

\subsection{RandomForest}

\textbf{RandomForest} è un classificatore basato su un insieme di alberi decisionali. Ogni albero produce una predizione e il modello finale combina i risultati. In genere è robusto e adatto a dataset complessi.

\subsection{NaiveBayes}

\textbf{NaiveBayes} è un classificatore probabilistico basato sul teorema di Bayes. Assume, in modo semplificato, che le feature siano indipendenti tra loro.

\begin{notaprof}
NaiveBayes è un classificatore molto veloce in Weka. Questa caratteristica può renderlo utile come baseline o come modello semplice da confrontare con algoritmi più complessi.
\end{notaprof}

\subsection{IBk}

\textbf{IBk} è l'implementazione Weka del metodo \emph{k-Nearest Neighbours}. È un classificatore \emph{instance-based}: una nuova istanza viene classificata in base alle classi delle \(k\) istanze più simili presenti nel training set.

\section{Validazione sperimentale}

La tecnica richiesta è la \textbf{10 times 10-fold cross-validation}.

Nella \textbf{10-fold cross-validation} semplice, il dataset viene diviso in 10 parti. A turno, 9 parti vengono usate per il training e 1 per il testing, fino a usare ciascuna fold come test set una volta.

Nella \textbf{10 times 10-fold cross-validation}, l'intero processo viene ripetuto 10 volte con partizioni differenti, ottenendo complessivamente 100 run.

Ripetere la cross-validation serve a ridurre l'effetto della casualità nella suddivisione delle fold e a ottenere una stima più stabile delle performance.

\begin{notaesame}
La 10 times 10-fold cross-validation è richiesta dalla Milestone 2. Tuttavia, nella Software Engineering è utile ragionare anche su split temporali, come ordered holdout o walk-forward, perché mischiare release passate e future può produrre una stima ottimistica dell'accuratezza.
\end{notaesame}

\section{Feature Selection}

La Milestone 2 richiede l'uso di filtri come tecnica di \textbf{feature selection}. La feature selection permette di selezionare solo le metriche più utili, eliminando attributi irrilevanti o ridondanti.

I principali vantaggi sono:

\begin{itemize}
    \item ridurre il rumore nel dataset;
    \item migliorare, in alcuni casi, le performance del modello;
    \item ridurre il costo di estrazione e misurazione delle metriche;
    \item semplificare l'interpretazione del modello.
\end{itemize}

\begin{notaprof}
Se dopo la feature selection l'accuratezza migliora, significa che alcune feature generavano confusione nel modello. Se invece l'accuratezza diminuisce leggermente, si può comunque ottenere un vantaggio pratico: usare meno metriche e ridurre il costo di raccolta dati.
\end{notaprof}

\begin{notaesame}
La feature selection non deve essere applicata sull'intero dataset prima della cross-validation. Se il filtro vede anche le fold di test, si introduce \textbf{data leakage} e la valutazione diventa falsata.
\end{notaesame}

In Weka, una soluzione corretta è usare il meta-classificatore:

\[
\texttt{AttributeSelectedClassifier}
\]

In questo modo la feature selection viene applicata separatamente sulle fold di training, senza usare informazioni provenienti dalla fold di test.

\section{Balancing}

Nei dataset di defect prediction la classe \textbf{buggy} è spesso minoritaria. Questo sbilanciamento può spingere il classificatore a predire quasi sempre la classe maggioritaria, cioè \texttt{non buggy}.

Per mitigare il problema si possono usare tecniche di \textbf{balancing}, come:

\begin{itemize}
    \item \textbf{undersampling}: riduzione della classe maggioritaria;
    \item \textbf{oversampling}: aumento della classe minoritaria;
    \item \textbf{SMOTE}: generazione di istanze sintetiche della classe minoritaria.
\end{itemize}

\begin{notaesame}
Il balancing deve essere applicato solo al \textbf{training set}, mai al test set. Bilanciare il test set significa alterare artificialmente la distribuzione reale dei dati su cui si valuta il modello.
\end{notaesame}

In Weka, il balancing può essere gestito tramite il meta-classificatore:

\[
\texttt{FilteredClassifier}
\]

\begin{notaesame}
Se si usano sia feature selection sia balancing, bisogna evitare che uno dei due passaggi contamini il test set. Le operazioni devono essere incapsulate nei meta-classificatori, in modo che vengano applicate solo alle fold di training durante la cross-validation.
\end{notaesame}

\begin{notaprof}
Quando si combinano feature selection e balancing, è importante ragionare sull'ordine delle operazioni. Il balancing altera artificialmente le istanze, quindi può influenzare la valutazione statistica delle feature se applicato nel momento sbagliato.
\end{notaprof}

\section{Metriche di valutazione}

\subsection{Precision}

La \textbf{Precision} misura quanti degli elementi predetti come buggy sono effettivamente buggy:

\[
Precision = \frac{TP}{TP + FP}
\]

Una Precision alta indica pochi falsi positivi.

\subsection{Recall}

La \textbf{Recall} misura quanti degli elementi realmente buggy sono stati trovati dal classificatore:

\[
Recall = \frac{TP}{TP + FN}
\]

Una Recall alta indica pochi falsi negativi.

\subsection{AUC}

L'\textbf{AUC}, cioè Area Under the ROC Curve, valuta la capacità del classificatore di distinguere tra classi positive e negative al variare della soglia decisionale.

\subsection{Kappa}

Il \textbf{Kappa di Cohen} confronta le prestazioni del classificatore con quelle ottenibili per caso. Un valore maggiore di zero indica che il modello sta facendo meglio di una classificazione casuale o banale.

\begin{notaprof}
Il classificatore ZeroR predice sempre la classe maggioritaria. In dataset sbilanciati può ottenere un'accuracy alta, ma un Kappa pari a zero, perché non sta realmente imparando un pattern predittivo.
\end{notaprof}

\subsection{NPofB20}

\begin{notaslide}
\textbf{NPofB20} è una metrica effort-aware normalizzata. Valuta quanti difetti si riescono a trovare ispezionando il 20\% delle LOC del sistema.
\end{notaslide}

La logica è ordinare le classi in base a una priorità di ispezione, tipicamente legata alla probabilità predetta di difettosità normalizzata rispetto alla dimensione della classe.

\begin{notaesame}
L'accuracy da sola non basta nei dataset sbilanciati. Un modello che predice sempre \texttt{non buggy} può sembrare accurato, ma non trovare nessuna classe difettosa. Per questo sono fondamentali metriche come Recall, Kappa e NPofB20.
\end{notaesame}

\section{Domande da rispondere nella Milestone 2}

La Milestone 2 richiede di rispondere a domande sperimentali come:

\begin{itemize}
    \item quale classificatore è più accurato degli altri?
    \item il miglior classificatore cambia in base al dataset?
    \item il miglior classificatore cambia in base al numero di release considerate?
    \item il miglior classificatore cambia in base alla metrica?
    \item quale classificatore è migliore per una determinata metrica?
\end{itemize}

Non esiste necessariamente un classificatore migliore in assoluto. Un modello può avere alta Recall ma bassa Precision, oppure buona AUC ma Kappa modesto.

\begin{notaesame}
Nel report è importante discutere i trade-off. Il docente non cerca solo il valore numerico migliore, ma la capacità di giustificare le scelte sperimentali e interpretare i risultati.
\end{notaesame}

\section{Weka e API}

\subsection{Uso della GUI}

La GUI di Weka, in particolare Explorer ed Experimenter, è utile per comprendere il workflow, testare rapidamente i classificatori e confrontare configurazioni diverse.

L'Experimenter permette di lanciare esperimenti ripetuti, come la 10 times 10-fold cross-validation, e salvare i risultati in una tabella esportabile.

\begin{notaprof}
Dal CSV prodotto dall'Experimenter può essere utile costruire una colonna descrittiva del modello, combinando classificatore e opzioni, ad esempio \texttt{RandomForest\_FS\_Balanced}.
\end{notaprof}

\subsection{Limite della GUI per NPofB20}

La GUI di Weka non calcola nativamente la \textbf{NPofB20}. Per calcolarla servono informazioni a livello di singola istanza, come:

\begin{itemize}
    \item ID dell'istanza;
    \item LOC o dimensione della classe;
    \item classe reale;
    \item probabilità predetta;
    \item ordinamento delle classi secondo la priorità effort-aware.
\end{itemize}

\begin{notaprof}
Le API Java di Weka o uno script esterno sono utili per stampare le predizioni istanza per istanza e calcolare automaticamente NPofB20.
\end{notaprof}

\section{Collegamento con il progetto}

L'input della Milestone 2 è il CSV finale prodotto nella Milestone 1. L'output è una tabella di confronto tra classificatori e configurazioni sperimentali.

Per ciascuno dei tre classificatori si possono confrontare diverse varianti:

\begin{itemize}
    \item modello base;
    \item modello con feature selection;
    \item modello con balancing;
    \item modello con feature selection e balancing.
\end{itemize}

Questo produce una tabella comparativa su più metriche, come Precision, Recall, AUC, Kappa e NPofB20.

\begin{notaesame}
L'obiettivo finale della Milestone 2 è identificare e giustificare un \textbf{best classifier}, cioè la combinazione di algoritmo e preprocessing più adatta al proprio dataset. Questo modello sarà poi usato nella Milestone 3.
\end{notaesame}

\section{Da ricordare}

\begin{notaesame}
Feature selection e balancing devono essere applicati senza contaminare il test set. In Weka è opportuno usare meta-classificatori come \texttt{AttributeSelectedClassifier} e \texttt{FilteredClassifier}.
\end{notaesame}

\begin{notaesame}
Il balancing va applicato solo al training set. Bilanciare il test set significa falsare la valutazione.
\end{notaesame}

\begin{notaesame}
La 10 times 10-fold cross-validation è richiesta dalla Milestone 2, ma bisogna essere consapevoli che nei dati software l'ordine temporale delle release ha significato.
\end{notaesame}

\begin{notaesame}
L'accuracy non è sufficiente per valutare modelli su dataset sbilanciati. Bisogna considerare Precision, Recall, AUC, Kappa e NPofB20.
\end{notaesame}

\begin{notaesame}
Il miglior classificatore può cambiare in base alla metrica. Per questo il risultato deve essere discusso, non solo riportato numericamente.
\end{notaesame}